\section{اشتقاق مكونات النموذج الديناميكي للذراع الروبوتية}

انطلاقاً من الشكل المصفوفي للنموذج الديناميكي العام للذراع الروبوتية ذات خمس درجات حرية، يمكن كتابة معادلات العزوم بالشكل التالي:

\begin{equation}
\resizebox{0.8\textwidth}{!}{$
\begin{bmatrix}
T_1 \\
T_2 \\
T_3 \\
T_4 \\
T_5
\end{bmatrix}
=
D_{ij}
\begin{bmatrix}
\ddot{\theta}_1 \\
\ddot{\theta}_2 \\
\ddot{\theta}_3 \\
\ddot{\theta}_4 \\
\ddot{\theta}_5
\end{bmatrix}
+
D_{ijj}
\begin{bmatrix}
\dot{\theta}_1^{2} \\
\dot{\theta}_2^{2} \\
\dot{\theta}_3^{2} \\
\dot{\theta}_4^{2} \\
\dot{\theta}_5^{2}
\end{bmatrix}
+
D_{ijk}
\begin{bmatrix}
\dot{\theta}_1\dot{\theta}_2 \\
\dot{\theta}_1\dot{\theta}_3 \\
\dot{\theta}_1\dot{\theta}_4 \\
\dot{\theta}_1\dot{\theta}_5 \\
\dot{\theta}_2\dot{\theta}_3 \\
\dot{\theta}_2\dot{\theta}_4 \\
\dot{\theta}_2\dot{\theta}_5 \\
\dot{\theta}_3\dot{\theta}_4 \\
\dot{\theta}_3\dot{\theta}_5 \\
\dot{\theta}_4\dot{\theta}_5
\end{bmatrix}
+
\begin{bmatrix}
D_1 \\
D_2 \\
D_3 \\
D_4 \\
D_5
\end{bmatrix}
$}
\end{equation}

حيث إن:

\begin{flushright}
\large
\begin{tabular}{@{}r@{}}
$D_{ij}$ : مصفوفة العطالة \\[0.1cm]

$D_{ijj}$ : مصفوفة القوى الطاردة المركزية \\[0.1cm]

$D_{ijk}$ : مصفوفة قوى كوريوليس \\[0.1cm]

$D_i$ : مصفوفة الجاذبية
\end{tabular}
\end{flushright}

تمثل العلاقة السابقة الشكل المصفوفي الكامل للنموذج الديناميكي للذراع الروبوتية، حيث يتألف النموذج من حدود العطالة، وحدود السرعات، وحدود الجاذبية. وسيتم في الأقسام التالية حساب هذه الحدود بشكل تفصيلي للوصول إلى معادلات العزوم الخاصة بمفاصل الذراع.










\subsection{حساب مصفوفة العطالة $D_{ij}$}

لحساب مصفوفة العطالة $D_{ij}$ لا بد أولاً من حساب تغير مصفوفات التحويل بالنسبة لمتحولات المفاصل، لأن عناصر هذه المصفوفة تعتمد على البنية الهندسية للذراع وعلى توزع الكتل وعزوم القصور الذاتي للوصلات.

تعطى مصفوفة المشتقات الجزئية لمصفوفة التحويل بالعلاقة:

\begin{equation}
U_{ij}
=
\frac{\partial T_i^0}{\partial q_j}
=
\frac{\partial \left(T_1 T_2 \cdots T_j \cdots T_i\right)}
{\partial q_j}
=
T_1 T_2 \cdots Q_j T_j \cdots T_i
\end{equation}

حيث إن $q_i$ تمثل متحولات المفاصل، وبما أن جميع مفاصل الذراع الروبوتية المدروسة هي مفاصل دورانية، فإن مصفوفة $Q$ الخاصة بالمفصل الدوراني تعطى بالشكل:

\begin{equation}
Q_{\text{revolute}} =
\begin{bmatrix}
0 & -1 & 0 & 0 \\
1 & 0  & 0 & 0 \\
0 & 0  & 0 & 0 \\
0 & 0  & 0 & 0
\end{bmatrix}
\end{equation}

وبالتالي يمكن حساب مصفوفات $U_{ij}$ اللازمة لمصفوفة العطالة كما يلي:

\begin{align}
U_{11} &= Q T_1
\\
U_{21} &= Q T_1T_2,
&
U_{22} &= T_1QT_2
\\
U_{31} &= Q T_1T_2T_3,
&
U_{32} &= T_1QT_2T_3,
&
U_{33} &= T_1T_2QT_3
\\
U_{41} &= Q T_1T_2T_3T_4,
&
U_{42} &= T_1QT_2T_3T_4
\\
U_{43} &= T_1T_2QT_3T_4,
&
U_{44} &= T_1T_2T_3QT_4
\\
U_{51} &= Q T_1T_2T_3T_4T_5,
&
U_{52} &= T_1QT_2T_3T_4T_5
\\
U_{53} &= T_1T_2QT_3T_4T_5,
&
U_{54} &= T_1T_2T_3QT_4T_5
\\
U_{55} &= T_1T_2T_3T_4QT_5
\end{align}

بعد حساب مصفوفات $U_{ij}$، يتم استخدام مصفوفة القصور الذاتي لكل وصلة من وصلات الذراع. وتكون مصفوفات القصور الذاتي للوصلات كما يلي:
\[
\resizebox{0.85\textwidth}{!}{$
\begin{array}{@{}c@{\hspace{1.2cm}}c@{}}

J_1 =
\begin{bmatrix}
2211 & 0 & 0 & -75 \\
0 & 3751 & 0 & -67 \\
0 & 0 & 1940 & 404.5 \\
-75 & -67 & 404.5 & 257.7
\end{bmatrix}
&
J_2 =
\begin{bmatrix}
377.88 & 0 & 0 & 0 \\
0 & 5632 & 0 & 690.8 \\
0 & 0 & 884.5 & 47.18 \\
0 & 690.5 & 47.18 & 168.5
\end{bmatrix}
\\[2cm]

J_3 =
\begin{bmatrix}
661.47 & 0 & 0 & 334 \\
0 & 670 & 0 & 50.5 \\
0 & 0 & 377.5 & -19 \\
334 & 50.5 & -19 & 90
\end{bmatrix}
&
J_4 =
\begin{bmatrix}
295 & 0 & 0 & 38.6 \\
0 & 1303.2 & 0 & -3.86 \\
0 & 0 & 8235 & 1323.9 \\
38.6 & -3.86 & 1323.9 & 193
\end{bmatrix}
\\[2cm]

\multicolumn{2}{c}{
J_5 =
\begin{bmatrix}
300 & 0 & 0 & 13.875 \\
0 & 2751 & 0 & 425.725 \\
0 & 0 & 376 & -297 \\
13.875 & 425.7 & -297 & 92.5
\end{bmatrix}
}

\end{array}
$}
\]





كما تم حساب كتلة ومركز كتلة كل وصلة بالنسبة إلى الإطار المحلي الخاص بها، مع إهمال القاعدة لأنها ثابتة ولا تدخل ضمن الوصلات المتحركة في النموذج الديناميكي. وكانت القيم كما يلي:

\[
\resizebox{0.7\textwidth}{!}{$
\begin{array}{@{}c@{\hspace{1.5cm}}c@{}}

m_1 = 257.7 \ \text{g},
\qquad
r_{c1} =
\begin{bmatrix}
-0.29 \\
-0.26 \\
1.57 \\
1
\end{bmatrix}
\ \text{cm}
&
m_2 = 168.5 \ \text{g},
\qquad
r_{c2} =
\begin{bmatrix}
0 \\
4.1 \\
0.28 \\
1
\end{bmatrix}
\ \text{cm}
\\[2cm]

m_3 = 90 \ \text{g},
\qquad
r_{c3} =
\begin{bmatrix}
3.71 \\
0.56 \\
-0.21 \\
1
\end{bmatrix}
\ \text{cm}
&
m_4 = 193 \ \text{g},
\qquad
r_{c4} =
\begin{bmatrix}
0.2 \\
-0.02 \\
6.86 \\
1
\end{bmatrix}
\ \text{cm}
\\[2cm]

\multicolumn{2}{c}{
m_5 = 92.5 \ \text{g},
\qquad
r_{c5} =
\begin{bmatrix}
0.15 \\
4.6 \\
-3.21 \\
1
\end{bmatrix}
\ \text{cm}
}

\end{array}
$}
\]

تُستخدم هذه القيم في حساب مركبات الجاذبية ومصفوفة العطالة، حيث يمثل كل متجه $r_{ci}$ مركز كتلة الوصلة $i$ بالنسبة إلى إطارها المحلي.
وتحسب عناصر مصفوفة العطالة بالاعتماد على العلاقة العامة:

\begin{equation}
D_{ij}
=
\sum_{p=\max(i,j)}^{5}
\operatorname{tr}
\left(
U_{pj}J_pU_{pi}^{T}
\right)
\end{equation}

وبتطبيق العلاقة السابقة على الذراع الروبوتية ذات خمس درجات حرية، نحصل على العناصر القطرية لمصفوفة العطالة كما يلي:

\begin{align*}
D_{11} &=
\operatorname{tr}(U_{11}J_1U_{11}^{T})
+
\operatorname{tr}(U_{21}J_2U_{21}^{T})
+
\operatorname{tr}(U_{31}J_3U_{31}^{T})
+
\operatorname{tr}(U_{41}J_4U_{41}^{T})
+
\operatorname{tr}(U_{51}J_5U_{51}^{T})
\\
D_{22} &=
\operatorname{tr}(U_{22}J_2U_{22}^{T})
+
\operatorname{tr}(U_{32}J_3U_{32}^{T})
+
\operatorname{tr}(U_{42}J_4U_{42}^{T})
+
\operatorname{tr}(U_{52}J_5U_{52}^{T})
\\
D_{33} &=
\operatorname{tr}(U_{33}J_3U_{33}^{T})
+
\operatorname{tr}(U_{43}J_4U_{43}^{T})
+
\operatorname{tr}(U_{53}J_5U_{53}^{T})
\\
D_{44} &=
\operatorname{tr}(U_{44}J_4U_{44}^{T})
+
\operatorname{tr}(U_{54}J_5U_{54}^{T})
\\
D_{55} &=
\operatorname{tr}(U_{55}J_5U_{55}^{T})
\end{align*}

أما العناصر غير القطرية فتحسب كما يلي:

\begin{align*}
D_{12} &=
\operatorname{tr}(U_{22}J_2U_{21}^{T})
+
\operatorname{tr}(U_{32}J_3U_{31}^{T})
+
\operatorname{tr}(U_{42}J_4U_{41}^{T})
+
\operatorname{tr}(U_{52}J_5U_{51}^{T})
\\
D_{13} &=
\operatorname{tr}(U_{33}J_3U_{31}^{T})
+
\operatorname{tr}(U_{43}J_4U_{41}^{T})
+
\operatorname{tr}(U_{53}J_5U_{51}^{T})
\\
D_{14} &=
\operatorname{tr}(U_{44}J_4U_{41}^{T})
+
\operatorname{tr}(U_{54}J_5U_{51}^{T})
\\
D_{15} &=
\operatorname{tr}(U_{55}J_5U_{51}^{T})
\\
D_{23} &=
\operatorname{tr}(U_{33}J_3U_{32}^{T})
+
\operatorname{tr}(U_{43}J_4U_{42}^{T})
+
\operatorname{tr}(U_{53}J_5U_{52}^{T})
\\
D_{24} &=
\operatorname{tr}(U_{44}J_4U_{42}^{T})
+
\operatorname{tr}(U_{54}J_5U_{52}^{T})
\\
D_{25} &=
\operatorname{tr}(U_{55}J_5U_{52}^{T})
\\
D_{34} &=
\operatorname{tr}(U_{44}J_4U_{43}^{T})
+
\operatorname{tr}(U_{54}J_5U_{53}^{T})
\\
D_{35} &=
\operatorname{tr}(U_{55}J_5U_{53}^{T})
\\
D_{45} &=
\operatorname{tr}(U_{55}J_5U_{54}^{T})
\end{align*}

وبما أن مصفوفة العطالة متناظرة، فإن:

\begin{equation*}
D_{ij}=D_{ji}
\end{equation*}

وبالتالي تكتب مصفوفة العطالة النهائية بالشكل:

\begin{equation*}
D(q)=
\begin{bmatrix}
D_{11} & D_{12} & D_{13} & D_{14} & D_{15} \\
D_{21} & D_{22} & D_{23} & D_{24} & D_{25} \\
D_{31} & D_{32} & D_{33} & D_{34} & D_{35} \\
D_{41} & D_{42} & D_{43} & D_{44} & D_{45} \\
D_{51} & D_{52} & D_{53} & D_{54} & D_{55}
\end{bmatrix}
\end{equation*}

تمثل هذه المصفوفة تأثير كتل الوصلات وعزوم قصورها الذاتي على حركة الذراع، وتعد من أهم مكونات النموذج الديناميكي، لأنها تحدد مقدار العزم اللازم لإنتاج تسارع معين عند كل مفصل.














\subsection{حساب متجه الجاذبية $D_g(q)$}

لحساب مركبات الجاذبية في النموذج الديناميكي، يجب التعبير عن شعاع الجاذبية بالنسبة إلى إطار القاعدة. وبما أن الحسابات في هذا النموذج تمت باستخدام السنتيمتر كوحدة للطول والغرام كوحدة للكتلة، فإن تسارع الجاذبية يكتب بالشكل:

\[
g_0 =
\begin{bmatrix}
0 \\
0 \\
-981 \\
0
\end{bmatrix}
\ \mathrm{cm/s^2}
\]

حيث يمثل $g_0$ شعاع الجاذبية المعبر عنه في إطار القاعدة. أما متجه مركز الكتلة $r_{cp}$ لكل وصلة، فهو معرف بالنسبة إلى الإطار المحلي للوصلة نفسها. وبما أن:

\[
U_{pi} =
\frac{\partial T_p^0}{\partial q_i}
\]

فإن الحد:

\[
U_{pi}r_{cp}
\]

يمثل تأثير موضع مركز كتلة الوصلة $p$ بالنسبة إلى المتغير المفصلي $q_i$، ومعبرًا عنه بالنسبة إلى إطار القاعدة. لذلك يجب استخدام شعاع الجاذبية $g_0$ المعبر عنه في إطار القاعدة نفسه عند حساب مركبات الجاذبية.

تحسب مركبة الجاذبية الموافقة للمفصل $i$ بالاعتماد على العلاقة العامة:

\[
D_i =
\sum_{p=i}^{5}
- m_p g_0^{T} U_{pi} r_{cp}
\]

حيث إن $m_p$ تمثل كتلة الوصلة $p$، و $r_{cp}$ يمثل متجه مركز كتلة الوصلة $p$ بالنسبة إلى إطارها المحلي، و $U_{pi}$ يمثل مشتقة مصفوفة التحويل الخاصة بالوصلة $p$ بالنسبة إلى متغير المفصل $i$.

وبتطبيق العلاقة السابقة على الذراع الروبوتية ذات خمس درجات حرية، نحصل على مركبات الجاذبية كما يلي:

\begin{align*}
D_1 &=
- m_1 \left(g_0^{T} U_{11} r_{c1}\right)
- m_2 \left(g_0^{T} U_{21} r_{c2}\right)
- m_3 \left(g_0^{T} U_{31} r_{c3}\right)
- m_4 \left(g_0^{T} U_{41} r_{c4}\right)
- m_5 \left(g_0^{T} U_{51} r_{c5}\right)
\\[0.2cm]
D_2 &=
- m_2 \left(g_0^{T} U_{22} r_{c2}\right)
- m_3 \left(g_0^{T} U_{32} r_{c3}\right)
- m_4 \left(g_0^{T} U_{42} r_{c4}\right)
- m_5 \left(g_0^{T} U_{52} r_{c5}\right)
\\[0.2cm]
D_3 &=
- m_3 \left(g_0^{T} U_{33} r_{c3}\right)
- m_4 \left(g_0^{T} U_{43} r_{c4}\right)
- m_5 \left(g_0^{T} U_{53} r_{c5}\right)
\\[0.2cm]
D_4 &=
- m_4 \left(g_0^{T} U_{44} r_{c4}\right)
- m_5 \left(g_0^{T} U_{54} r_{c5}\right)
\\[0.2cm]
D_5 &=
- m_5 \left(g_0^{T} U_{55} r_{c5}\right)
\end{align*}

وبذلك يمكن كتابة متجه الجاذبية للنظام بالشكل:

\[
D_g(q)=
\begin{bmatrix}
D_1 \\
D_2 \\
D_3 \\
D_4 \\
D_5
\end{bmatrix}
\]

يمثل هذا المتجه العزوم الناتجة عن تأثير أوزان الوصلات على مفاصل الذراع. ويعد حد الجاذبية من الحدود المهمة في هذا النموذج، لأن الذراع الروبوتية تعمل ضد تأثير الجاذبية أثناء رفع الأجسام أو المحافظة على وضعية ثابتة.






















\subsection{إضافة عزم القصور الذاتي للمشغلات}

إضافة إلى قصور الوصلات، يجب أخذ عزم القصور الذاتي للمشغلات بعين الاعتبار، لأن كل محرك يضيف مقاومة إضافية للتسارع عند المفصل الموافق. وقد تم اعتماد القيمة التقريبية التالية لعزم القصور الذاتي للمشغل:

\[
I_{act}=4\times10^{-6} \ \text{kg.m}^2
\]

وبما أن الحسابات السابقة تمت بوحدات الغرام والسنتيمتر، يتم تحويل هذه القيمة إلى وحدة $g.cm^2$ كما يلي:

\[
1 \ \text{kg.m}^2 = 10^{7} \ \text{g.cm}^2
\]

ومن ثم:

\[
I_{act}
=
4\times10^{-6} \times 10^{7}
=
40 \ \text{g.cm}^2
\]

وبما أن الذراع تحتوي على خمسة مفاصل، يمكن تمثيل عزوم القصور الذاتي للمشغلات بمصفوفة قطرية بالشكل:

\[
I_{act} =
\begin{bmatrix}
40 & 0 & 0 & 0 & 0 \\
0 & 40 & 0 & 0 & 0 \\
0 & 0 & 40 & 0 & 0 \\
0 & 0 & 0 & 40 & 0 \\
0 & 0 & 0 & 0 & 40
\end{bmatrix}
\ \text{g.cm}^2
\]

وبالتالي يصبح الحد العطالي الكلي في النموذج الديناميكي مكوّناً من مصفوفة عطالة الوصلات $D(q)$ بالإضافة إلى مصفوفة قصور المشغلات $I_{act}$، أي:

\[
D_{\text{total}}(q)
=
D(q)+I_{act}
\]

وتستخدم هذه المصفوفة الكلية في حساب العزوم المطلوبة لتسريع مفاصل الذراع الروبوتية.



\subsection{إهمال قوى كوريوليس والقوى النابذة المركزية}

في النموذج الديناميكي العام للذراع الروبوتية تظهر حدود مرتبطة بسرعات المفاصل، وهي حدود قوى كوريوليس والقوى النابذة المركزية. وتكتب هذه الحدود ضمن النموذج العام بالشكل:

\[
D_{ijj}
\begin{bmatrix}
\dot{\theta}_1^2 \\
\dot{\theta}_2^2 \\
\dot{\theta}_3^2 \\
\dot{\theta}_4^2 \\
\dot{\theta}_5^2
\end{bmatrix}
+
D_{ijk}
\begin{bmatrix}
\dot{\theta}_1\dot{\theta}_2 \\
\dot{\theta}_1\dot{\theta}_3 \\
\dot{\theta}_1\dot{\theta}_4 \\
\dot{\theta}_1\dot{\theta}_5 \\
\dot{\theta}_2\dot{\theta}_3 \\
\dot{\theta}_2\dot{\theta}_4 \\
\dot{\theta}_2\dot{\theta}_5 \\
\dot{\theta}_3\dot{\theta}_4 \\
\dot{\theta}_3\dot{\theta}_5 \\
\dot{\theta}_4\dot{\theta}_5
\end{bmatrix}
\]

حيث تمثل المصفوفة $D_{ijj}$ معاملات القوى النابذة المركزية الناتجة عن مربعات السرعات الزاوية، بينما تمثل المصفوفة $D_{ijk}$ معاملات قوى كوريوليس الناتجة عن جداءات السرعات الزاوية بين المفاصل المختلفة.

في هذا المشروع تم إهمال هذين الحدين في النموذج التحليلي المستخدم لحساب العزوم، وذلك اعتماداً على طبيعة الذراع الروبوتية المدروسة وطبيعة المهمة المطلوبة منها. فالذراع المستخدمة صغيرة الحجم، وتتعامل مع أوزان خفيفة نسبياً، كما أن حركتها أثناء نقل الأجسام لا تتم بسرعات عالية، وإنما تكون حركة بطيئة ومحدودة بهدف الوصول إلى الوضعيات المطلوبة بشكل آمن ومستقر.

تعتمد قوى كوريوليس والقوى النابذة المركزية بشكل مباشر على السرعات الزاوية للمفاصل. فالقوى النابذة المركزية تتناسب مع مربعات السرعات الزاوية مثل $\dot{\theta}_i^2$، بينما تتناسب قوى كوريوليس مع جداءات السرعات الزاوية مثل $\dot{\theta}_i\dot{\theta}_j$. لذلك، عندما تكون السرعات الزاوية صغيرة، تصبح هذه الحدود صغيرة جداً مقارنة بحدود الجاذبية والعطالة.

كما أن الجزء الأكبر من العزم المطلوب في هذه الذراع ينتج عن مقاومة تأثير الجاذبية، خاصة في المفاصل التي تحمل وزن عدة وصلات مثل مفصل الكتف ومفصل الكوع. لذلك يكون تأثير متجه الجاذبية $D_i$ أكثر وضوحاً من تأثير قوى كوريوليس والقوى النابذة المركزية، وخصوصاً عند دراسة الوضعيات الساكنة أو الحركة البطيئة.

إضافة إلى ذلك، فإن الهدف من النموذج الديناميكي في هذا الفصل هو بناء نموذج تحليلي مبسط يساعد على فهم السلوك الديناميكي العام للذراع وتقدير العزوم المطلوبة عند المفاصل، وليس الوصول إلى نموذج ديناميكي كامل عالي الدقة. وبما أن إدخال حدود كوريوليس والقوى النابذة المركزية يؤدي إلى زيادة كبيرة في طول وتعقيد المعادلات، فقد تم إهمالها للحصول على نموذج أبسط وأكثر ملاءمة للتحليل.

وبناءً على ذلك يمكن اعتماد التقريب التالي:

\[
D_{ijj}
\begin{bmatrix}
\dot{\theta}_1^2 \\
\dot{\theta}_2^2 \\
\dot{\theta}_3^2 \\
\dot{\theta}_4^2 \\
\dot{\theta}_5^2
\end{bmatrix}
\approx
0
\]

\[
D_{ijk}
\begin{bmatrix}
\dot{\theta}_1\dot{\theta}_2 \\
\dot{\theta}_1\dot{\theta}_3 \\
\dot{\theta}_1\dot{\theta}_4 \\
\dot{\theta}_1\dot{\theta}_5 \\
\dot{\theta}_2\dot{\theta}_3 \\
\dot{\theta}_2\dot{\theta}_4 \\
\dot{\theta}_2\dot{\theta}_5 \\
\dot{\theta}_3\dot{\theta}_4 \\
\dot{\theta}_3\dot{\theta}_5 \\
\dot{\theta}_4\dot{\theta}_5
\end{bmatrix}
\approx
0
\]

وعليه يصبح النموذج الديناميكي المستخدم في الحساب كما يلي:

\[
T =
\left(D(q)+I_{act}\right)\ddot{q}
+
D_g(q)
\]

أي أن العزوم المطلوبة للمفاصل سيتم حسابها بالاعتماد على مصفوفة العطالة، ومصفوفة قصور المشغلات، ومتجه الجاذبية فقط.

تجدر الإشارة إلى أن إهمال قوى كوريوليس والقوى النابذة المركزية لا يعني أنها غير موجودة فيزيائياً، وإنما يعني أن تأثيرها في شروط عمل هذا المشروع صغير مقارنة ببقية الحدود الديناميكية. أما في حال تشغيل الذراع بسرعات عالية، أو حمل أوزان كبيرة، أو تنفيذ مسارات سريعة ومعقدة، فيجب إعادة إدراج هذه الحدود للحصول على نموذج ديناميكي أكثر دقة.